\par\Needspace{12\baselineskip}\bigskip
\pdfbookmark[2]{Exposition, Section deuxième}{pallas-II-section2}
\begin{center}
{\large\bfseries Section deuxième.\par}
\medskip
\textit{Variations instantanées des élémens, produites par les perturbations.}
\end{center}
\pdfbookmark[3]{Exposition, 7}{pallas-II-article7}
\begin{center}7.\end{center}
\noindent\hspace*{\parindent}Le mouvement elliptique, que nous venons d’expliquer, suppose que les attractions mutuelles du soleil et de la planète soient les seules forces qui agissent. S’il y en a encore d’autres forces qui viennent en considération, le mouvement de la planète autour du soleil devient plus compliqué. Mais quelqu’il soit, on voit que, si ces forces étaient anéanties tout d’un coup, la planète se mouvrait, dès ce moment, dans une section conique, dont les élémens seraient fonctions des quantités qui pour cet instant déterminent le lieu et la vitesse acquise de la planète, c. à d. des six quantités \(x,y,z,\dfrac{\mathrm{d}x}{\mathrm{d}t}\mathpunct{\raisebox{.45ex}{,}}\dfrac{\mathrm{d}y}{\mathrm{d}t}\mathpunct{\raisebox{.45ex}{,}}\dfrac{\mathrm{d}z}{\mathrm{d}t}\mathpunct{\cdot}\) Nous nous bornons au cas où cette section conique est toujours une ellipse, \SourceLeaf{456}et nous la nommerons \textls[80]{ellipse osculatrice}. On voit donc que toutes les équations finies entre les élémens et les quantités \(x,y,z,r,\dfrac{\mathrm{d}x}{\mathrm{d}t}\mathpunct{\raisebox{.45ex}{,}}\dfrac{\mathrm{d}y}{\mathrm{d}t}\mathpunct{\raisebox{.45ex}{,}}\dfrac{\mathrm{d}z}{\mathrm{d}t}\mathpunct{\raisebox{.45ex}{,}}\dfrac{\mathrm{d}r}{\mathrm{d}t}\mathpunct{\raisebox{.45ex}{,}}\) que nous avons développées pour le cas du mouvement purement elliptique, ou plutôt toutes celles qui ont lieu pour ce cas, seront également vraies pour l’ellipse osculatrice. Mais comme les forces perturbatrices ne cessent point d’agir continuellement, l’ellipse osculatrice variera d’un instant à l’autre: ses élémens seront des quantités variables. Cependant si l’on parvient à pouvoir assigner pour chaque instant les valeurs numériques de tous les élémens de l’ellipse osculatrice, le problème de la détermination du mouvement sera résolu: on se servira des mêmes formules comme dans le cas du mouvement non troublé pour calculer non seulement le lieu de la planète, mais aussi sa vitesse et la direction de son mouvement. Cette recherche se réduit à deux problèmes: d’abord nous déterminerons les variations instantanées que prennent les différens élémens par l’action des forces perturbatrices: après cela il faudra intégrer ces expressions différentielles, pour obtenir pour chaque instant les valeurs des élémens variables.
